\section{Compound AI Systems 的经验规律}

\subsection{核心模式对比}

\begin{table}[H]
\centering
\caption{Compound AI Systems 的构建模式}
\begin{tabular}{p{0.2\textwidth}p{0.26\textwidth}p{0.22\textwidth}p{0.24\textwidth}}
\toprule
模式 & 组成 & 优势 & 工程挑战 \\
\midrule
RAG & 检索 + 1 次 LLM & 实现简单 & 缺乏链式推理能力 \\
ReAct & 推理 → 工具 → 观察 循环 & 适合多步骤任务 & 需要精细的控制逻辑 \\
DSPy & 声明式结构 + 优化器 & 可编程化、可迁移 & 搜索空间大、调优耗时 \\
Autonomous Agent & 多 Agent 协作 + MCP & 自动化并行工作流 & 需要强监控与安全护栏 \\
\bottomrule
\end{tabular}
\end{table}

\begin{knowledgebox}{经验规律}
复合 AI 系统的设计遵循三个准则：1) 组件需要明确签名以便组合；2) 结构必须被 metric / log 检验；3) 优化过程要可回放，否则系统容易 drift。
\end{knowledgebox}

\subsection{签名与组合}

Signature 定义了模块的输入/输出形式，DSPy 通过 signature 检查避免不合法的组合。每次 pipeline 扩展时，都应插入类型检查器并自动生成交互测试，防止下游模块收到奇异值。

\subsection{模式重用}

复合系统常见模式（如 retriever → reasoner、planner → executor）可以封装成 reusable module。一旦签名一致，就能在不同 pipeline 中循环利用，大幅度减少手工 prompt 复制的开销。

\begin{knowledgebox}{复用的前提是签名一致}
再强的优化器也无法保护不匹配的组合。如果一个 ``state → action'' 模块被插错到 ``question → answer'' 的 pipeline，就会在 Evidence Matrix 中不断生成 invalid token。签名检查器要在 compile 阶段就抛出错误。
\end{knowledgebox}

\begin{warningbox}{不要把 pipeline 当成黑箱}
一个包含检索、推理、工具、校验的 pipeline 中任何一步出错都可能导致错误传播，所以在每个模块边界设置 guard rails（断言、类型检查、fallback）比优化 prompt 更重要。
\end{warningbox}

\subsection{本章小结}

Compound AI Systems 的可靠性来自清晰接口、可组合性检验与可观测的优化流程，而不是仅仅依赖更强的 LLM。

